\documentclass[11pt]{amsart}
\usepackage[T1]{fontenc}
\usepackage[utf8]{inputenc}
\usepackage{lmodern,microtype}
\usepackage[margin=27mm]{geometry}
\usepackage{amsmath,amssymb,amsthm,mathtools}
\usepackage{xcolor}
\usepackage[colorlinks=true,linkcolor=blue!50!black,citecolor=blue!50!black,urlcolor=blue!50!black]{hyperref}
\usepackage{xurl}
\setlength{\emergencystretch}{2em}

\newcommand{\workbenchurl}{https://github.com/KokunoYumeto/grothendieck-questions-workbench}
\newcommand{\recordurl}{https://github.com/KokunoYumeto/mathematics-workbenches}

\newtheorem{theorem}{Theorem}
\newtheorem{thm}{Theorem}[section]
\newtheorem{lemma}[thm]{Lemma}
\newtheorem{proposition}[thm]{Proposition}
\newtheorem{corollary}[thm]{Corollary}
\theoremstyle{remark}
\newtheorem{remark}[thm]{Remark}

\newcommand{\Z}{\mathbb Z}
\newcommand{\Q}{\mathbb Q}
\newcommand{\F}{\mathbb F}
\DeclareMathOperator{\Sec}{Sec}
\DeclareMathOperator{\Inn}{Inn}
\DeclareMathOperator{\Aut}{Aut}
\DeclareMathOperator{\divi}{div}

\title[Prime-to-$p$ sections separate rational points]{Rational points of curves in characteristic $p$ are separated by the prime-to-$p$ fundamental group, on a question of Grothendieck}
\author{}
\date{23 November 2026}

\begin{document}

\begin{abstract}
In a manuscript now in the Montpellier archive, Grothendieck observed that over a field of characteristic $p$, and with prime-to-$p$ fundamental groups, the abelian part of the fundamental group no longer separates rational points, already for a punctured elliptic curve, because the prime-to-$p$ Kummer map does not see $p$-power torsion; he asked whether the whole prime-to-$p$ fundamental group still separates them. Sa\"idi and Tamagawa proved that for a hyperbolic curve over a finite field the decomposition groups in the geometrically prime-to-$p$ arithmetic fundamental group separate closed points outside a finite exceptional set. We prove that for every hyperbolic curve $X$ over a field $K$ finitely generated over $\F_p$, the map from $X(K)$ to conjugacy classes of sections of the geometrically prime-to-$p$ arithmetic fundamental group is injective: the exceptional set disappears for rational points. The proof uses no conjecture; its inputs are SGA~1, the Weil conjectures for curves and the finiteness of Sa\"idi and Tamagawa's exceptional set. The results and their proofs were produced by GPT-6 Astra (ChatGPT Pro mode) and checked by Claude Opus 5.5 at effort max (an independent re-proof of the first two steps, and line-by-line checks of the others), which also wrote this text; no human has checked them.
\end{abstract}

\maketitle

\section{The question and the result}

\subsection{Grothendieck's question}
Let $K$ be a field, $X/K$ a smooth geometrically connected curve and $\overline K$ a separable closure. A rational point $P\in X(K)$ gives a section $s_P$ of $\pi_1(X)\to G_K=\mathrm{Gal}(\overline K/K)$, well defined up to conjugation by $\pi_1(X_{\overline K})$, and the first question of the arithmetic of fundamental groups is whether $P\mapsto[s_P]$ is injective. Over finitely generated fields of characteristic $0$ it is, by Kummer theory and the Mordell--Weil theorem: two distinct points have distinct classes already in the abelianization.

In characteristic $p$ one must replace the geometric fundamental group by a quotient that behaves well in families, and the natural one is the maximal prime-to-$p$ quotient. In a manuscript of the Montpellier archive \cite[cote 149, ms.~pp.~17--18]{Archive}, Grothendieck notes that with prime-to-$p$ fundamental groups the abelian version of injectivity fails, already for an elliptic curve with a point removed: the prime-to-$p$ Kummer map of the Jacobian kills its $p$-primary torsion. He asks whether the full prime-to-$p$ version may still hold. He reduces it to a covering statement: if $P,Q$ have the same class, then for every connected prime-to-$p$ Galois cover and every point $P'$ above $P$ there is a point $Q'$ above $Q$ with $Q'-P'$ of $p$-power order in the generalized Jacobian of the cover. He writes that one might be able to deduce $P=Q$ from this.

\subsection{The theorem}
Let $p$ be a prime, $K$ a field finitely generated over $\F_p$, and $X/K$ a smooth geometrically connected curve which is hyperbolic: $2g-2+r>0$, where $g$ is the genus of its smooth compactification and $r$ the number of geometric points at infinity. Put
\[
 \Delta_X=\pi_1(X_{\overline K})^{(p')},\qquad
 \Pi_X=\pi_1(X)\big/\ker\bigl(\pi_1(X_{\overline K})\to\Delta_X\bigr),
\]
where $(p')$ denotes the maximal prime-to-$p$ quotient, so that $1\to\Delta_X\to\Pi_X\to G_K\to1$ is exact and the arithmetic quotient is the whole of $G_K$.

\begin{theorem}\label{thm:main}
The map
\[
 X(K)\longrightarrow \Sec(\Pi_X\to G_K)/\Delta_X,\qquad P\longmapsto[s_P],
\]
is injective.
\end{theorem}

In Grothendieck's covering form: if $P,Q\in X(K)$ satisfy his condition for all prime-to-$p$ covers, then $P=Q$. For proper curves over a finite field the two forms are equivalent by Kummer theory, since the prime-to-$p$ Kummer kernel of a finite abelian group is its $p$-primary part. The hypothesis that $X$ be hyperbolic excludes $\mathbb A^1$, whose geometric prime-to-$p$ fundamental group is trivial, so that the map is constant; for $\mathbb G_m$ the Kummer argument of Section~\ref{sec:open} gives injectivity as well.

\subsection{Why it matters}
For the full fundamental group, closed points of hyperbolic curves over finite fields are separated by their decomposition groups, by Tamagawa for affine curves and by Mochizuki for proper ones; this is the starting point of anabelian geometry over finite fields, as Sa\"idi and Tamagawa recall \cite[Introduction, p.~3]{ST09}. For the prime-to-$p$ quotient, the natural one for families and for the questions in Grothendieck's manuscript, Sa\"idi and Tamagawa proved that the decomposition-group map is injective outside a finite exceptional set \cite[Definitions 1.3, 1.5, 1.10 and Proposition 1.8(vi)]{ST09}, and wrote that they could only prove it ``almost injective'' \cite[p.~3]{ST09}. In their later refinement for geometrically pro-$\Sigma$ groups the exceptional sets could not even be shown to be finite \cite[Introduction, pp.~4--5]{ST18}. Theorem~\ref{thm:main} shows that for rational points no exceptional set is needed, over finite fields and over every finitely generated field of characteristic $p$, and so answers Grothendieck's question for rational points of curves.

\subsection{Credit}
The question and the covering formulation are Grothendieck's \cite[cote 149, ms.~pp.~17--18]{Archive}. The finite exceptional set used in Step~4 is Sa\"idi and Tamagawa's \cite{ST09}, and so is the commensurator lemma used in Step~5 \cite[Lemma 1.7]{ST09}. The comparison theorems are those of SGA~1 \cite{SGA1}, and the weights of Frobenius on $H^1$ of a curve are those of the Weil conjectures for curves \cite[(1.5.1) and Th\'eor\`eme 1.6]{WeilI}. What is new here is the argument that removes the exceptional set: the transport of the cusp through the configuration space of two points, the comparison of point counts it gives, and the Kummer argument with the function whose divisor is $p^e(P-Q)$.

\subsection{Context}
Injectivity for the full fundamental group over finitely generated fields is classical (by Kummer theory and Mordell--Weil in characteristic $0$, and by Tamagawa and Mochizuki over finite fields); the open part of the section conjecture is surjectivity. In characteristic $p$, Grothendieck's own remark shows why the prime-to-$p$ version is harder: the abelian information is blind to $p$-power torsion of Jacobians. We know of no earlier proof of Theorem~\ref{thm:main}; we have not been able to consult Stix's monograph on rational points and fundamental groups, which treats the injectivity part of the section conjecture.

\subsection{Scope}
Proved: Theorem~\ref{thm:main} for rational points, over every field finitely generated over $\F_p$, for proper and for affine hyperbolic curves. Not proved: separation of all closed points, i.e.\ the removal of Sa\"idi and Tamagawa's exceptional set for decomposition groups of points that are not rational over the base; the analogous statement for geometrically pro-$\Sigma$ quotients with $\Sigma$ not cofinite; any statement on surjectivity.

\subsection{Proof overview}
Suppose $P\ne Q$ have conjugate sections in $\Pi_C$, for a proper curve $C$ of genus at least two over a finite field $k$. The difference $P-Q$ is then $p$-power torsion in the Jacobian, which is exactly what prime-to-$p$ Kummer theory cannot see; so the argument must use the nonabelian structure. It does so in three moves.

First, the configuration space $C\times C-\Delta$ fibres over $C$ with fibre $C$ minus a point. Fix a prime $r\ne p$. After passing to the quotient that is pro-$r$ along the punctures, the fibre sequence stays exact (Step~1), and the decomposition group of the diagonal carries the full $\Z_r(1)$ of the puncture (Step~2). A conjugator between $s_P$ and $s_Q$ then moves the fibre over $P$ to the fibre over $Q$ while fixing the compactified curve, giving an isomorphism $\beta$ between the cuspidally pro-$r$ fundamental groups of $C-P$ and $C-Q$, over the identity of $\Pi_C$.

Second, $\beta$ preserves the integral $r$-adic cohomology of the compactifications of corresponding covers, with Frobenius. By the trace formula, corresponding covers have the same number of points over every finite extension (Step~3). A compactness argument then shows that $\beta$ sends decomposition groups of rational points to decomposition groups of rational points; outside Sa\"idi and Tamagawa's finite exceptional set, the point is the same one (Step~4).

Third, the function $f$ with divisor $p^e(P-Q)$ is compared with $\beta$ on the cuspidally central quotients. Its Kummer class differs from $p^e$ times a class built from $\beta$ by a constant, because $H^1$ of the curve with $\Z_r(1)$ coefficients has no Frobenius invariants (weights). Evaluated at a good point $R$, the class built from $\beta$ vanishes, so $f(R)$ is, up to a fixed constant $a$, of order prime to $r$. But $f$ takes every value, among them $a\zeta$ for a root of unity $\zeta$ of order a power of $r$ (Step~5). The general case reduces to this one by specialization and by looking at the compactification (Section~\ref{sec:open}).

\section{The finite-field case}\label{sec:finite}

In this section $k$ is a finite field of characteristic $p$ with $q$ elements, $C/k$ is a proper smooth geometrically connected curve of genus $g\ge2$, and $P,Q\in C(k)$. We write $\Pi_C$, $\Delta_C$ as in Section~1, and we fix a prime $r\ne p$. For a profinite group $K$ let $N(K)$ be the intersection of the open normal subgroups of $K$ of $r$-power index, so that $K^{(r)}=K/N(K)$; $N(K)$ is characteristic in $K$.

\begin{thm}\label{thm:finite}
If $s_P$ and $s_Q$ are conjugate by an element of $\Delta_C$, then $P=Q$.
\end{thm}

The rest of this section proves Theorem~\ref{thm:finite}; suppose $P\ne Q$ and $s_Q=\gamma s_P\gamma^{-1}$ with $\gamma\in\Delta_C$.

\subsection{Step 1: the configuration fibration stays exact}
Let $F=\Delta_{C-P}$ and $E=\Delta_{C\times C-\mathrm{diag}}$ be the geometric prime-to-$p$ fundamental groups, $K_F=\ker(F\to\Delta_C)$ and $K_E=\ker(E\to\Delta_C\times\Delta_C)$, and put $F_c=F/N(K_F)$ and $E_c=E/N(K_E)$; the second projection $C\times C-\mathrm{diag}\to C$ has fibre $C-P$ over $P$.

\begin{lemma}\label{lem:free}
Let $1\to L\to H\to B\to1$ be an exact sequence of groups with $L$ free of finite rank $d\ge2$, such that $H$ acts on $L/[L,L]L^r$ through a finite $r$-group. Then $L^{(r)}\to H^{(r)}$ is injective.
\end{lemma}

\begin{proof}
Conjugation gives an action of $H$ on $L^{(r)}$, and $\Aut(L^{(r)})=\varprojlim\Aut(L^{(r)}/U)$ over the open characteristic subgroups $U$. For a finite $r$-group $Q$ with Frattini subgroup $\Phi(Q)$, the kernel of $\Aut(Q)\to\Aut(Q/\Phi(Q))$ is an $r$-group: by Burnside's basis theorem it acts freely on the $|\Phi(Q)|^d$ tuples lifting a fixed basis of $Q/\Phi(Q)$. For $Q=L^{(r)}/U$, the group $Q/\Phi(Q)$ is a quotient of $L/[L,L]L^r$, on which $H$ acts through an $r$-group; so $H$ acts on each $Q$ through an $r$-group, and the closure $A$ of the image of $H$ in $\Aut(L^{(r)})$ is pro-$r$. Hence $H\to A$ factors through $H^{(r)}$. On $L$ the action is inner, so $L^{(r)}\to H^{(r)}\to A$ is $x\mapsto\Inn(x)$, whose kernel is the centre of $L^{(r)}$, trivial for a free pro-$r$ group of rank at least two. So $L^{(r)}\to H^{(r)}$ is injective.
\end{proof}

\begin{proposition}\label{prop:exact}
The sequence $1\to F_c\to E_c\to\Delta_C\to1$ induced by the second projection is exact, and the same holds for the arithmetic quotients $\Pi_{C-P}/N(K_F)$ and $\Pi_{C\times C-\mathrm{diag}}/N(K_E)$ over $G_k$.
\end{proposition}

\begin{proof}
First let $C$ be a compact Riemann surface of genus $g\ge2$, and use topological fundamental groups $\Lambda=\pi_1(C-P)$, $\Gamma=\pi_1(C\times C-\mathrm{diag})$, $B_0=\pi_1(C)$. The second projection is a locally trivial fibration with fibre $C-P$ (Fadell--Neuwirth), and $\pi_2(C)=0$, so $1\to\Lambda\to\Gamma\to B_0\to1$ is exact. Prime-to-$p$ completion is right exact; $N(K_F)$ maps into $N(K_E)$; and $E_c\to\Delta_C$ is onto with kernel the image of $F_c$. It remains to show that $F_c\to E_c$ is injective, that is, that every finite quotient $F_c\to\mathcal Q$ is dominated by a finite quotient of $E_c$.

Let $R\subset\mathcal Q$ be the image of $K_F$, an $r$-group, and $A=\mathcal Q/R$, a quotient of $F/K_F=\Delta_C$. Let $a\colon C'\to C$ be the corresponding Galois cover with group $A$, and $Y=(C'\times C')-\bigcup_{\sigma\in A}\{(\sigma x,x)\}$; these graphs are disjoint because $A$ acts freely, and $Y\to C\times C-\mathrm{diag}$ is Galois with group $A\times A$, so $H=\pi_1(Y)$ is normal in $\Gamma$ with $\Gamma/H=A\times A$. Since $\Lambda$ maps to $B_0\times\{1\}$, we have $\Lambda\cap H=\ker(\Lambda\to A)=L$, the fundamental group of the fibre $V'=C'-a^{-1}(P)$ of the second projection $Y\to C'$; $L$ is free of rank $2g(C')+|A|-1\ge4$. The monodromy of $\pi_1(C')$ on $H_1(V',\F_r)$ is trivial on the quotient $H_1(C',\F_r)$, because the compactified family is the product $C'\times C'\to C'$, and trivial on the classes of the punctures, which are not permuted; so it is unipotent with square zero, an elementary abelian $r$-group, and Lemma~\ref{lem:free} applies to $1\to L\to H\to\pi_1(C')\to1$. The image of $L$ in $\mathcal Q$ is $R$, an $r$-group, so there is a normal subgroup $T\subset H$ of $r$-power index with $T\cap L\subseteq\ker(L\to R)$. Let $T'=\bigcap_{\gamma\in\Gamma/H}\gamma T\gamma^{-1}$; it is normal in $\Gamma$, $H/T'$ is an $r$-group, and $\Gamma/T'$ is an extension of $A\times A$ by $H/T'$, of order prime to $p$. So $\Gamma\to\Gamma/T'$ factors through $E$, and through $E_c$, because $K_E$ maps into the $r$-group $H/T'$. Finally $\Lambda\cap T'\subseteq\Lambda\cap H\cap T=L\cap T\subseteq\ker(L\to R)\subseteq\ker(\Lambda\to\mathcal Q)$, using $\Lambda\cap H=L$. So $\Lambda\to\mathcal Q$ factors through $\Lambda\to\Gamma/T'$, which proves injectivity.

For the curve over $k$, lift $C$ and $P$ to $W(k)$ (SGA~1, III, Corollaire 7.4; the point lifts by smoothness). The prime-to-$p$ geometric fundamental groups of $C-P$, of $C\times C-\mathrm{diag}$ (a smooth proper surface minus a smooth divisor, along which every prime-to-$p$ cover is tame) and of $C$ specialize isomorphically and compatibly \cite[XIII, Corollaire 2.9, \S2.10, Lemme 2.11]{SGA1}. The subgroups $K_F,K_E,N(\cdot)$ are defined group-theoretically, so the sequence is exact for the \'etale groups. As $N(K_F)$ and $N(K_E)$ are characteristic in normal subgroups of the arithmetic groups, the arithmetic quotients exist and are extensions of the same $G_k$.
\end{proof}

\subsection{Step 2: equal sections transport the cusp}
Let $\mathcal G=\Pi_{C\times C-\mathrm{diag}}/N(K_E)$ and $G_P=\Pi_{C-P}/N(K_F)$; by Proposition~\ref{prop:exact}, $G_P$ is the inverse image of $s_P(G_k)$ under the second projection $\mathcal G\to\Pi_C$, and similarly for $Q$. Let $\mathcal D\subset\mathcal G$ be a decomposition group of the diagonal.

\begin{lemma}\label{lem:diag}
$\mathcal D$ is an extension $1\to\Z_r(1)\to\mathcal D\to\Pi_C\to1$, mapping diagonally to $\Pi_C\times_{G_k}\Pi_C$, and $\mathcal D\cap G_P$ is a decomposition group of the cusp of $C-P$.
\end{lemma}

\begin{proof}
Prime-to-$p$ covers of $C\times C-\mathrm{diag}$ are tame along the smooth divisor $\mathrm{diag}\cong C$, so by Abhyankar's lemma on Kummer charts \cite[XIII, Proposition 5.2]{SGA1} the decomposition group of the diagonal is an extension of $\pi_1(\mathrm{diag})$, whose image is $\Pi_C$ embedded diagonally, by the image of tame inertia, which in $\mathcal G$ is a quotient of $\Z_r(1)$. It is the whole of $\Z_r(1)$: with $\Lambda=\langle a_i,b_i\rangle$, the loop around the puncture is $\prod_i[a_i,b_i]$, and the map $a_1\mapsto1+e_{12}$, $b_1\mapsto1+e_{23}$, other generators $\mapsto1$, into the pro-$r$ Heisenberg group of upper unitriangular $3\times3$ matrices over $\Z_r$, sends it to $[1+e_{12},1+e_{23}]=1+e_{13}$, of infinite order; this map factors through $F_c$, and by Proposition~\ref{prop:exact} through $E_c$. The diagonal meets each fibre transversally in one point, so the cusp inertia of the fibre is the diagonal inertia, and the cusp group of the fibre over $P$ is $\mathcal D\cap G_P$.
\end{proof}

Lift $\gamma$ to $\tilde\gamma\in\mathcal D$; then both projections of $\tilde\gamma$ to $\Pi_C$ equal $\gamma$, and conjugation by $\tilde\gamma$ maps $G_P$ onto $G_Q$ and the cusp group $\mathcal D\cap G_P$ onto $\mathcal D\cap G_Q$, acting on the first projection by $\Inn(\gamma)$. Since $F_c\to\Delta_C$ is onto, there is $\delta$ in the geometric part of $G_Q$ whose first projection is $\gamma^{-1}$. Conjugation by $\delta\tilde\gamma$ is an isomorphism
\[
 \beta\colon G_P\xrightarrow{\ \sim\ }G_Q
\]
over the identity of $\Pi_C$ (through the first projections) and over $G_k$, which preserves the conjugacy classes of cusp groups.

\subsection{Step 3: corresponding covers have the same number of points}

\begin{lemma}\label{lem:rquot}
Let $K\triangleleft F$ be closed and $N=N(K)$. Every open subgroup $I\le F$ containing $N$ satisfies $I^{(r)}=(I/N)^{(r)}$.
\end{lemma}

\begin{proof}
Let $I/J$ be a finite $r$-quotient, $I_0$ the core of $I$ in $F$, and $M$ the core of $J\cap I_0$ in $F$. Then $I_0/M$ is an $r$-group, and the image of $K$ in $F/I_0$ is an $r$-group because $N\subseteq I_0$; so the image of $K$ in $F/M$ is an $r$-group, hence $N\subseteq M\subseteq J$.
\end{proof}

Let $I$ be an open subgroup of $G_P$ and $V_I\to(C-P)_{k_I}$ the corresponding cover, $k_I$ its field of constants, and $X_I$ its smooth compactification. The geometric fundamental group of $V_I$ has the same pro-$r$ completion as its image in the geometric prime-to-$p$ group, and by Lemma~\ref{lem:rquot} this is the pro-$r$ completion of the geometric part of $I$. Killing the inertia of the cusps gives the geometric pro-$r$ fundamental group of $X_I$, with its Frobenius action. Since $\beta$ is an isomorphism over $G_k$ preserving cusp classes, it identifies $H^1(X_I\otimes\overline k,\Z_r)$ with $H^1(X_{\beta(I)}\otimes\overline k,\Z_r)$ as modules over the Frobenius of $k_I$. By the Lefschetz trace formula for curves, $X_I$ and $X_{\beta(I)}$ have the same number of points over every finite extension of $k_I$. The cusps of $V_I$ over $P$ and their degrees are read off from the double cosets of the cusp groups, which $\beta$ preserves. Subtracting, $V_I$ and $V_{\beta(I)}$ have the same number of points over every finite extension of $k_I$.

\subsection{Step 4: $\beta$ fixes the points outside a finite set}
Let $L/k$ be finite, $R\in(C-P)(L)$, $D_R\subset (G_P)_L$ a decomposition group, and $S=\beta(D_R)$. For every open $H\supseteq S$ of $(G_Q)_L$, the cover corresponding to $\beta^{-1}(H)\supseteq D_R$ has an $L$-point above $R$, so by Step~3 the cover $V_H$ of $C-Q$ corresponding to $H$ has an $L$-point. The sets $V_H(L)$, $H\supseteq S$, are finite and nonempty and form an inverse system over the directed set of such $H$, so their limit is nonempty. A compatible family gives $T\in(C-Q)(L)$ and a decomposition group $D_T$ contained in every $H\supseteq S$, hence in $S$; both project isomorphically onto $G_L$, so $D_T=S$. Thus $\beta$ sends point sections to point sections, and, being over the identity of $\Pi_C$, it sends the proper decomposition group of $R$ to that of $T$.

By Sa\"idi and Tamagawa there is a finite set $E_C$ of closed points of $C$, for the cofinite set of primes different from $p$, outside which the decomposition group in $\Pi_C$, even up to commensurability, determines the point \cite[Definitions 1.3, 1.5, 1.10, Proposition 1.8(vi)]{ST09}; it serves for every constant extension \cite[Remark 1.2]{ST09}. Hence $T=R$ whenever $R\notin E_C\cup\{P,Q\}$.

\subsection{Step 5: the Kummer contradiction}
Equal sections have equal images in the abelianization, so the class of $P-Q$ lies in the kernel of the prime-to-$p$ Kummer map of the finite group $J_C(k)$, which is its $p$-primary part. Choose $e\ge0$ and $f\in k(C)^\times$ with $\divi f=p^e(P-Q)$.

Let $M=\Z_r(1)$, and let $E_P$, $E_Q$ be the quotients of $G_P$, $G_Q$ that are central along the cusp: $1\to M\to E_P\to\Pi_C\to1$, with extension class in $H^2(\Delta_C,M)=\Z_r$ the class of the point (the first Chern class of $\mathcal O(P)$, by the presentation of the punctured surface, SGA~1, XIII, Corollaire 2.12); by Lemma~\ref{lem:diag} the kernel is all of $M$. The isomorphism $\beta$ induces $\bar\beta\colon E_P\to E_Q$ over $\Pi_C$, acting on $M$ by a unit $u$; as both extension classes are the class of a point, $u=1$. On $E_P\times_{\Pi_C}E_Q$, $\chi(x,y)=\bar\beta(x)y^{-1}$ is an $M$-valued crossed homomorphism. Pulled back to the cuspidally central quotient of the fundamental group of $C-\{P,Q\}$, it has residues $1$ at $P$ and $-1$ at $Q$, while the Kummer class $\kappa(f)\in H^1(\,\cdot\,,M)$ has residues $p^e$ and $-p^e$. So $\kappa(f)-p^e\chi$ is unramified and lies in $H^1(\Pi_C,M)$. By inflation--restriction, and since $G_k$ has cohomological dimension one, $H^1(\Pi_C,M)=H^1(G_k,M)$: the Frobenius eigenvalues on $\mathrm{Hom}(T_rJ_C,\Z_r(1))\otimes\Q_r$ have absolute value $\sqrt q$ \cite[Th\'eor\`eme 1.6]{WeilI}, so there are no invariants, and the module is torsion-free. By Kummer theory over $k$ there is $a\in k^\times$ with
\[
 \kappa(f)=p^e\chi+\kappa(a).
\]
Let $R\notin E_C\cup\{P,Q\}$ be a closed point and $L=k(R)$. By Step~4, $\bar\beta(D_R^P)$ is conjugate to $D_R^Q$; the conjugator projects to the normalizer of $D_R$ in $\Pi_{C_L}$, which is $D_R$ itself \cite[Lemma 1.7]{ST09}, so after multiplying by an element of $D_R^Q$ it lies in $M$. Then $\chi$ restricted to the decomposition group of $R$ is the coboundary of an element of $M$, and restricting the displayed identity to it gives $\kappa_L(f(R)/a)=0$: the element $f(R)/a$ has order prime to $r$ in $L^\times$.

Since $P\ne Q$, $f$ is not constant, so $f\colon C\to\mathbb P^1$ is onto. The points of $E_C\cup\{P,Q\}$ have finitely many values under $f$; choose a root of unity $\zeta\in\overline k$ of order a power of $r$, $\zeta\ne1$, with $a\zeta$ different from these values, and a closed point $R$ with $f(R)=a\zeta$. Then $f(R)/a=\zeta$ has nontrivial $r$-power order in $k(R)^\times$, a contradiction. This proves Theorem~\ref{thm:finite}.

\section{Finitely generated fields and affine curves}\label{sec:open}

\subsection{Proper curves of genus at least two}
Let $K$ be finitely generated over $\F_p$, $C/K$ proper of genus $g\ge2$, and $P\ne Q$ in $C(K)$ with conjugate sections. Spread $C$, $P$, $Q$ out over a smooth connected scheme $S$ of finite type over $\F_p$ with function field $K$, shrinking $S$ so that $P$ and $Q$ extend to disjoint sections. Let $s\in S$ be a closed point, $A$ the henselization of $\mathcal O_{S,s}$ and $L$ its fraction field. Smooth proper specialization identifies the geometric prime-to-$p$ fundamental groups of $C_{\overline L}$ and $C_{\overline{k(s)}}$ \cite[X, Th\'eor\`eme 3.8, Corollaire 3.9]{SGA1}, and $G_L\to\pi_1(\mathrm{Spec}\,A)=G_{k(s)}$ is onto because $A$ is a normal henselian domain. Restricting the conjugacy to $G_L$ and pushing it to $\Pi_{C_A}=\Pi_{C_{k(s)}}$, the sections of the distinct points $P_s,Q_s\in C(k(s))$ become conjugate, against Theorem~\ref{thm:finite}.

\subsection{Affine curves}
Let $X\subset C$ be affine hyperbolic. After a finite separable extension of $K$ we may assume the points at infinity rational; a finite purely inseparable extension changes neither the \'etale fundamental groups nor the Galois group \cite[IX, Th\'eor\`eme 4.10]{SGA1}, and injectivity over an extension implies it over $K$. If $g\ge2$, the surjection $\Pi_X\to\Pi_C$ carries conjugate sections to conjugate sections, and the proper case applies.

If $g=0$, choose three points at infinity as $0,1,\infty$; the coordinate maps $X$ into $\mathbb G_m$, and conjugate sections give equal prime-to-$p$ Kummer classes of $z(P)$ and $z(Q)$ in $K^\times$. The kernel of the prime-to-$p$ Kummer map of $K^\times$ is trivial: modulo its finite group of constants, of order prime to $p$, $K^\times$ embeds into the free abelian group of divisors of a normal model, so an element that is an $n$-th power for every $n$ prime to $p$ is $1$. Hence $z(P)=z(Q)$ and $P=Q$.

If $g=1$, let $E=C$, with origin a rational point at infinity, let $S$ be the set of points at infinity, and let $P\ne Q$ in $X(K)$ have conjugate sections. The generalized Jacobian $G$ of $E$ with reduced modulus $S$ is an extension of $E$ by a split torus $T$. The group $E(K)$ is finitely generated (Lang--N\'eron; the field of constants of $K$ is finite), and $T(K)$ is a finite group of order prime to $p$ times a free abelian group. So an element $x$ of $G(K)$ divisible by every integer prime to $p$ is torsion of $p$-power order: its image in $E(K)$ is then torsion of $p$-power order; let $p^A$ be the exponent of $E(K)[p^\infty]$ and $N'$ that of the prime-to-$p$ torsion of $E(K)$. For $n$ prime to $p$ write $x=nN'y$; the image of $y$ in $E(K)$ is torsion, so $w=p^AN'y$ lies in $T(K)$ and $p^Ax=nw$. Thus $p^Ax\in T(K)$ is divisible in $T(K)$ by every integer prime to $p$, and by the argument of the case $g=0$, applied to each coordinate of $T(K)=(K^\times)^{\dim T}$, $p^Ax=0$. Hence the class of $Q-P$ in $G(K)$, which lies in the kernel of the prime-to-$p$ Kummer map, has some order $p^a$.

Choose $m$ prime to $p$ with $m^2>p^a$, and let $X'=E-[m]^{-1}(S)$, a Galois cover of $X$ with group $E[m]$, of degree $m^2$ prime to $p$; after a finite extension $K'$ of $K$ the points of $E[m]$ and of $[m]^{-1}(S)$ are rational. The section $s_Q|_{G_{K'}}=\gamma s_P|_{G_{K'}}\gamma^{-1}$ lands in $\Pi_{X'}$, which is normal, so $Q$ lifts to a point $Q'\in X'(K')$; writing $\gamma$ as an element of $\Delta_{X'}$ times a lift of a deck transformation $\tau$, the sections of $Q'$ and of the point $P'=\tau(P'_0)$ above $P$, for a lift $P'_0$ of $P$, are conjugate by $\Delta_{X'}$. So the class of $Q'-P'$ has $p$-power order in the generalized Jacobian $G'$ of $E$ with modulus $[m]^{-1}(S)$. Its image in $E$ has $m(Q'-P')=Q-P$, of order dividing $p^a$, and $m$ is prime to $p$; so it has order $p^b$ with $b\le a$, and since the torus of $G'$ has no nonzero $p$-primary points, $p^b$ is also its order in $G'$. Hence there is a function $h$ on $E$ with $\divi h=p^b(Q'-P')$ and $h=1$ at every point of $[m]^{-1}(S)$, a set of $m^2|S|\ge m^2$ points. As $P'\ne Q'$, $h-1$ is a nonzero function whose only pole is $P'$, of order $p^b<m^2$; so it has at most $p^b$ zeros, a contradiction. This proves Theorem~\ref{thm:main}.

\section{What remains open}
Sa\"idi and Tamagawa's exceptional set concerns all closed points, through their decomposition groups; Theorem~\ref{thm:main} removes it for points rational over the base field, where sections are available, but not for closed points of higher degree, whose decomposition groups are not sections. For geometrically pro-$\Sigma$ quotients with $\Sigma$ another cofinite set of primes, the same argument may apply where the configuration step and a finite exceptional set are available for $\Sigma$; we have not checked this; for $\Sigma$ not cofinite even finiteness of the exceptional sets is open \cite[pp.~4--5]{ST18}. Higher-dimensional versions, and surjectivity, are not addressed.

\section*{Note on method}
The results and their proofs were produced by GPT-6 Astra (ChatGPT Pro mode). Claude Opus 5.5 at effort max re-proved Steps~1 and~2 independently from the statements; in a separate check, Claude Opus 5.5 at effort max checked Steps~3--5 and Section~\ref{sec:open} line by line, corrected the genus-one argument of Section~\ref{sec:open} (the lifted points are paired by a deck transformation; no translation is needed), and wrote this text. No human has checked this note. Our text is dedicated to the public domain under CC0~1.0. The workbench is \expandafter\url\expandafter{\workbenchurl}; the permanent records are at \expandafter\url\expandafter{\recordurl}.

\begin{thebibliography}{9}
\bibitem{Archive} A.~Grothendieck, manuscripts, Archives Grothendieck, Universit\'e de Montpellier, cote 149, ms.~pp.~17--18 (scans 35--36 of the digitized cote). \url{https://grothendieck.umontpellier.fr/archives-grothendieck/}.
\bibitem{SGA1} A.~Grothendieck, \emph{Rev\^etements \'etales et groupe fondamental} (SGA~1), Lecture Notes in Math.\ 224, Springer, 1971: III 7.3--7.4; IX 4.10; X 3.8--3.9; XIII 2.9--2.12 and 5.2.
\bibitem{ST09} M.~Sa\"idi and A.~Tamagawa, \emph{A prime-to-$p$ version of Grothendieck's anabelian conjecture for hyperbolic curves over finite fields of characteristic $p>0$}, Publ.\ RIMS \textbf{45} (2009), 135--186; arXiv:0801.1452: Introduction, p.~3; Remark 1.2; Definitions 1.3, 1.5, 1.10; Lemma 1.7; Proposition 1.8.
\bibitem{ST18} M.~Sa\"idi and A.~Tamagawa, \emph{A refined version of Grothendieck's anabelian conjecture for hyperbolic curves over finite fields}, J.\ Algebraic Geom.\ \textbf{27} (2018), 383--448; arXiv:1708.08773, Introduction, pp.~4--5.
\bibitem{WeilI} P.~Deligne, \emph{La conjecture de Weil.~I}, Publ.\ Math.\ IH\'ES \textbf{43} (1974), 273--307: (1.5.1) and Th\'eor\`eme 1.6.
\end{thebibliography}

\end{document}
